\documentclass[]{style}

\newcommand{\printbibsectionkeyword}[3]{
  \begin{refsection}
    \newrefcontext[sorting=chronological]%
    \nocite{*}
    \printbibliography[type={#1}, title={#2}, keyword={#3}, heading=bibheading]
  \end{refsection}
}

\usepackage{fontawesome}
\usepackage{array}
\usepackage{hyperref}
\usepackage{tabularx}
\usepackage{booktabs}
\usepackage[ddmmyyyy]{datetime}


\newcolumntype{Y}{>{\centering\arraybackslash}X}


\addbibresource{bibliography.bib}
\addbibresource{misc.bib}


\begin{document}

\header{@firstname}{@lastname}{@qualification}


\begin{aside}
\section{{\color{white}.}}
\includegraphics[width=4cm]{me}
\section{Contact}
\begin{tabular}{m{1cm}m{3cm}}
    \faicon{globe}& \href{@url_website}{@website} \\
    \faicon{mobile} & \href{@url_phone}{@phone} \\
    \faicon{at} & \href{@url_email}{@email} \\
    \faicon{paper-plane} & \href{@url_telegram}{@telegram} \\
    \faicon{github} & \href{@url_github}{@github} \\
\end{tabular}
\section{Where}
\begin{tabular}{m{1cm}m{3cm}}
    \faicon{map} & @city \\
    \faicon{map-pin} & @street \\
    \faicon{building} & @building
\end{tabular}
\section{Languages}
\begin{tabular}{m{1cm}m{3cm}}
  \faicon{comments-o} & Italian \\
  \faicon{comments-o} & English 
\end{tabular}
\section{Programming}
\begin{tabular}{m{1cm}m{3cm}}
  \faicon{code} & Python \\
  \faicon{code} & Javascript \\
  \faicon{code} & C++ \\
  \faicon{code} & Java 
\end{tabular}
\end{aside}

Abstract goes here

\section{metrics}

\renewcommand{\arraystretch}{1.5}
\begin{table}[h]
  \centering
  \begin{tabularx}{\textwidth}{c X X X}
      & \centering Google Scholar & \centering Scopus \tabularnewline
      Articles & \centering 1 & \centering 2 \tabularnewline
      Citations & \centering 1 & \centering 2 \tabularnewline
      H-Index & \centering 1 & \centering 2 \tabularnewline
  \end{tabularx}
  \label{tab:metrics}
\end{table}
\renewcommand{\arraystretch}{1}

\section{education}

\begin{entrylist}

\entry
{2007--2013}
{Diploma {\normalfont Liceo Scientifico}}
{Liceo Scientifico G. Marconi, Grosseto (Italy)}

\entry
{2013--2016}
{Bachelor Degree {\normalfont in Computer Engineering}}
{University of Pisa, Pisa (Italy)}

\entry
{2016--2018}
{Master Degree {\normalfont in Computer Engineering}}
{University of Pisa, Pisa (Italy)}

\entry
{2018--2021}
{PhD {\normalfont in Information Engineering}}
{University of Pisa, Pisa (Italy)}

\end{entrylist}

\section{work}
\begin{entrylist}
\entry
{2021--2023}
{Research Fellow {\normalfont for the SERICA Project}}
{University of Pisa, Pisa (Italy)}

\entry
{2023--Now}
{Assistant Professor}
{University of Pisa, Pisa (Italy)}

\end{entrylist}


\section{academic projects}
\begin{entrylist}
  
\entry
{PRA 2022\_101}
{Decision Support Systems for Territorial Networks in Ecosystem Services Management}
{University of Pisa}
{Artificial intelligence technologies for monitoring water resources using environmental and satellite data.}

\entry
{Secure B2C}
{POR FESR 2014-2020}
{Italy}
{Development of AI prototypes for enterprise network security and hardware/software architecture for smart electronic payment devices. Published and presented at an international conference..}

\entry
{MUR-FISR 2019}
{SERICA}
{University of Pisa}
{Design of a digital library using AI for text analysis and interactive maps. Contributed significantly to AI components and supported the development of these systems. Published and presented at an international conference.}

\entry
{H2020}
{EXPERIENCE}
{Centro E. Piaggio}
{Integration of virtual reality and psychological components for virtual simulations. Developed a neural network for fNIRS signal regression.}

\entry
{Ki-Foot}
{POR FESR 2014-2020}
{Italy}
{Prototype for gait monitoring based on sensor integration in footwear. Collaborated on undergraduate theses.}

\entry
{PRA 2018\_81}
{University Research Project}
{University of Pisa}
{AI-based solutions for specific research problems. Published and presented at an international conference.}

\entry
{SCIADRO}
{Environmental Security and Protection with Drones}
{Italy}
{Use of drone swarms for detecting gas leaks and early fires. Developments based on computational stigmergy and bio-inspired algorithms. Published and presented at an international conference.}
  
\end{entrylist}

\section{publications}

\printbibsection{article}{Articles in international peer-reviewed journals} 
\printbibsection{inproceedings}{International peer-reviewed conferences} 

\section{roles}

\printbibsectionkeyword{misc}{For International Journals}{journalrole}
\printbibsectionkeyword{misc}{For International Conferences}{conferencerole}

\section{open-source projects}

\printbibsectionkeyword{misc}{Personal open-source projects}{project}

@legalstart

\newgeometry{left=2.5cm,top=2cm,right=2.5cm,bottom=2.5cm,nohead,nofoot}
@legalnotice
\vspace{2em}
\begin{table}[htbp]
  \centering
  \begin{tabularx}{\textwidth}{ Y  Y  Y }
      Pisa,  &     & @firstname @lastname \\
      \today &     &                      \\
             &     & \includegraphics[scale=1.5]{@signature_file} \\
  \end{tabularx}
\end{table}

@legalend

\end{document}
